\documentclass[11pt]{article}
\usepackage[review]{acl}
\usepackage{times}
\usepackage{latexsym}
\usepackage[T1]{fontenc}
\usepackage[utf8]{inputenc}
\usepackage{microtype}
\usepackage{graphicx}
\usepackage{booktabs}

\title{Does Self-Supervised Domain Fine-Tuning Help a Small Medical Retriever?\\
An Audit of a Published Encoder and a Preregistered Re-Training}

\author{Anonymous submission}

\begin{document}
\maketitle

\begin{abstract}
A small public project released sentence encoders fine-tuned from all-MiniLM-L6-v2 for a medical question-answering chatbot and reported a large gain on an internal similarity task. We could not reproduce that figure because the training code and data splits are missing, so we test the released checkpoints directly on four external retrieval benchmarks (SciFact, NFCorpus, TREC-COVID, and a MedQuAD question-to-answer task) against base, open general-purpose, and biomedical baselines. The released checkpoints fall far below their own base model on all four, and a diagnostic shows their embedding space is strongly compressed. We then re-train under a preregistered protocol on open MedlinePlus text with document-level splits. A first recipe falls below the base model; a second, stronger recipe (hard negatives and general-data replay, selected on held-out data) matches the base model within noise on three benchmarks and stays slightly below it on the fourth. Neither beats strong off-the-shelf encoders. In a small retrieval-augmented question-answering test with a gold-abstract corpus, retrieval raises a 7B generator's yes/no accuracy by about 0.17, and the published encoder's weaker retrieval gives lower accuracy than the others. Every weight tensor in the project's published LoRA adapters is NaN, so they cannot be used. We do not evaluate faithfulness, hallucination, or clinical safety.
\end{abstract}

\section{Introduction}
Domain-adapted embedding models are widely suggested as a cheap way to improve retrieval-augmented systems in specialised fields such as medicine. Public projects often report gains on a task built from the same corpus the model was trained on. Such a number says little about whether retrieval improves for new questions or new documents.

We audit one such project, HealthMate-AI: a chatbot whose retriever is all-MiniLM-L6-v2 \citep{reimers2019sbert,wang2020minilm} fine-tuned on text from a copyrighted medical encyclopedia. Its documentation reports a Spearman correlation of 0.8039 against 0.6795 for the base model. Our contributions are modest and specific:
\begin{itemize}
\item We document that the headline figure cannot be traced to code, and that a course report from the same project gives a different figure for the same ensemble (Section~\ref{sec:audit}).
\item We evaluate the released checkpoints zero-shot on four external retrieval sets against the systems in Table~\ref{tab:ndcg}, with bootstrap confidence intervals and Holm-corrected paired tests.
\item We re-train the same base model with two recipes under protocols fixed before the runs (the second added after the first failed), using only openly licensed MedlinePlus text, and report the results even though they are negative or null.
\end{itemize}
We make no claim about answer faithfulness, hallucination, or clinical use.

\section{Related Work}
BEIR \citep{thakur2021beir} provides zero-shot retrieval evaluation across domains, including the biomedical sets we use \citep{wadden2020scifact,boteva2016nfcorpus,voorhees2020trec}. General-purpose small encoders such as BGE \citep{xiao2023cpack}, E5 \citep{wang2022e5} and GTE \citep{li2023gte} are strong zero-shot baselines. MedCPT \citep{jin2023medcpt} is a biomedical retriever trained on PubMed search logs. Unsupervised domain adaptation of dense retrievers has dedicated methods, for example generative pseudo-labelling \citep{wang2022gpl} and denoising auto-encoder pre-training \citep{wang2021tsdae}; we do not compare against them. MedEIR \citep{selvadurai2025medeir} is one example of a specialised medical embedding model; we do not evaluate it. Networks trained sequentially on new data can forget earlier tasks (catastrophic forgetting) \citep{kirkpatrick2017ewc}, a possible but untested explanation for our results; averaging the weights of fine-tuned models can improve accuracy and robustness \citep{wortsman2022soups}. Our setting differs from these works in being an audit of a specific released artefact, not a new method. We did not perform a systematic literature survey, and we make no priority claim.

\section{The Audited Project and What Can Be Verified}\label{sec:audit}
The project's repository and model cards describe a 3-fold cross-validated fine-tune whose fold checkpoints were weight-averaged, and a 2-fold variant. Training pairs, according to a course report from the same project, were retrieved text chunks paired with reference answers and filtered by lexical overlap. The repository does not contain the training or evaluation scripts, the question files, or the splits; we therefore cannot verify how the folds were built or whether test pairs overlap training pairs. The same course report states two different ensemble results (0.8039 and 0.7552) with different base scores. We treat neither as a finding. Whatever the correct value is, it measures agreement on an in-corpus similarity task, not retrieval on unseen documents.

\section{Experimental Setup}
\paragraph{Benchmarks.} SciFact (300 test queries, 5,183 documents), NFCorpus (323 queries, 3,633 documents), TREC-COVID (50 queries, 171,332 documents), all as distributed with BEIR. We also build a MedQuAD \citep{benabacha2019medquad} task: each question must retrieve its own answer from 14,798 unique answers (15,395 questions) drawn from eight NIH sources. We exclude the MedlinePlus Health Topics subset because it overlaps our training text, and the three subsets whose answers were withheld for copyright. Questions whose answer is a verbatim duplicate treat every copy as relevant.

\paragraph{Systems.} BM25 (Okapi, \texttt{rank-bm25}); the base encoder; BGE-small, E5-small, GTE-small (used with their documented prefixes); MedCPT (query and article encoders); the two published HealthMate checkpoints; our re-trained models; and reciprocal-rank fusion (RRF, $k{=}60$) \citep{cormack2009rrf} of BM25 with BGE-small and with our model. BM25 is a standard baseline \citep{robertson2009bm25}.

\paragraph{Re-training.} We fine-tune all-MiniLM-L6-v2 on the English MedlinePlus health topics (public-domain text, 1,012 topics; 919 train, 93 validation by a fixed hash of the topic id). Each topic summary is cut into chunks of at most 120 words; pairs are (title, chunk) and (meta description, chunk), 6,666 in total, with in-batch negatives \citep{henderson2017smartreply}. We chose learning rate and epochs on validation only (a grid of three by three at seed 13); the best setting was the edge of the grid (learning rate $5\times10^{-5}$, 3 epochs). The validation task is nearly saturated (the base model scores 0.997 MRR@10), so it is weakly informative. Final models use seeds 13, 42 and 2024 and a uniform weight average of the three. An answer-level near-duplicate check between MedQuAD answers and training chunks (TF-IDF cosine) found 4 of 14,798 answers above 0.8.

\paragraph{Stronger recipe (follow-up).} After the first recipe underperformed, we preregistered a second study (Amendment 6, written before its runs, though all earlier results were known). Recipe R1 adds one hard negative per pair: a chunk from a different topic ranked 5 to 30 by the base model. Recipe R2 adds to R1 an equal number of general question--answer pairs from Natural Questions as replay (its licence was not confirmed; the pairs were used locally and not redistributed). We searched learning rate $\in\{5\times10^{-6},10^{-5},2\times10^{-5}\}$ and 1 or 2 epochs for each recipe (12 runs, seed 13), selecting on a new development split: MedQuAD questions from the GARD and GHR sources (10,796 queries). Evaluation of the final models uses the test split: the six remaining MedQuAD sources (4,599 queries) plus SciFact, NFCorpus and TREC-COVID. All twelve settings scored below the base model on the development split (0.607 to 0.659 against 0.693); the best was R2 with learning rate $5\times10^{-6}$ and one epoch, which we trained with seeds 13, 42, 2024 and averaged.

\paragraph{Metrics and statistics.} nDCG@10 (primary), Recall@10, Recall@100 and MRR@10, with 95\% percentile bootstrap intervals over queries (10,000 resamples). Paired bootstrap differences against the base model, with Holm correction across all comparisons reported. Which external baseline to compare our model against was chosen after seeing results, so that comparison is exploratory. Nothing was tuned on any benchmark. The preregistration and its seven dated amendments are in the repository; the zero-shot baseline runs and the runs of the published checkpoints were performed before the final protocol was frozen, which we state because it limits what ``preregistered'' covers.

\paragraph{Compute and software.} All retrieval and training runs used one Apple M4 Max laptop (38.6 GB unified memory, MPS backend, Python 3.12.13) with torch 2.14.1, sentence-transformers 6.1.0, transformers 5.18.0 and rank-bm25 0.2.2; the generator study used Mistral-7B-Instruct-v0.2 in float16. Each fine-tuning run took 2 to 4 minutes; zero-shot evaluation of one system took from seconds (SciFact, NFCorpus) to 38 minutes (the slowest system on TREC-COVID). Training used the sentence-transformers trainer with MultipleNegativesRankingLoss, batch size 32, 10\% linear warm-up, the library's default AdamW settings, a maximum sequence length of 256 tokens. Retrieval results are deterministic given the model; the training seeds are 13, 42 and 2024, and we report each seed and their average.

\paragraph{Data checks.} All benchmark and training text is public health or scientific text; we did not find or look for personal data, and we did not exhaustively inspect it for offensive content. We use each resource for retrieval evaluation or fine-tuning, consistent with its public release; the HealthMate checkpoints were released for medical retrieval and we test them on that use.

\section{Results}
Table~\ref{tab:ndcg} and Figure~\ref{fig:ndcg} give nDCG@10 for every system.

\begin{table*}[t]\centering\small
\begin{tabular}{lcccc}
\toprule
System & SciFact & NFCorpus & TREC-COVID & MedQuAD \\
\midrule
BM25 & 0.652 [0.61, 0.70] & 0.306 [0.27, 0.34] & 0.578 [0.50, 0.66] & 0.502 [0.50, 0.51] \\
all-MiniLM-L6-v2 (base) & 0.645 [0.60, 0.69] & 0.316 [0.28, 0.35] & 0.472 [0.39, 0.56] & 0.642 [0.64, 0.65] \\
BGE-small & 0.713 [0.67, 0.75] & 0.343 [0.31, 0.38] & 0.756 [0.70, 0.81] & 0.723 [0.72, 0.73] \\
E5-small & 0.687 [0.64, 0.73] & 0.325 [0.29, 0.36] & 0.744 [0.68, 0.81] & 0.713 [0.71, 0.72] \\
GTE-small & 0.727 [0.68, 0.77] & 0.349 [0.31, 0.38] & 0.665 [0.59, 0.74] & 0.725 [0.72, 0.73] \\
MedCPT & 0.735 [0.70, 0.77] & 0.353 [0.32, 0.39] & 0.594 [0.52, 0.67] & 0.535 [0.53, 0.54] \\
BM25+BGE (RRF) & 0.709 [0.66, 0.75] & 0.344 [0.31, 0.38] & 0.741 [0.68, 0.80] & 0.626 [0.62, 0.63] \\
Published HealthMate 3-fold & 0.249 [0.21, 0.29] & 0.069 [0.05, 0.08] & 0.219 [0.16, 0.29] & 0.180 [0.17, 0.18] \\
Published HealthMate best-2 & 0.250 [0.21, 0.29] & 0.068 [0.05, 0.08] & 0.233 [0.17, 0.30] & 0.177 [0.17, 0.18] \\
Ours, seed-averaged & 0.580 [0.54, 0.63] & 0.293 [0.26, 0.33] & 0.443 [0.36, 0.53] & 0.547 [0.54, 0.55] \\
\quad single seeds (min--max) & 0.576--0.583 & 0.284--0.295 & 0.437--0.441 & 0.540--0.544 \\
BM25+Ours (RRF) & 0.656 [0.61, 0.70] & 0.326 [0.29, 0.36] & 0.601 [0.52, 0.68] & 0.561 [0.56, 0.57] \\
\bottomrule
\end{tabular}

\caption{nDCG@10 [95\% bootstrap CI over queries]. Zero-shot; no system was tuned on these sets. ``Ours'' is trained on MedlinePlus only.}\label{tab:ndcg}
\end{table*}

\begin{figure*}[t]\centering
\includegraphics[width=\textwidth]{fig_ndcg.pdf}
\caption{nDCG@10 with 95\% confidence intervals for selected systems.}\label{fig:ndcg}
\end{figure*}

\paragraph{The published checkpoints.} On every set they score far below the base model (for example SciFact 0.249 against 0.645; paired differences and Holm-adjusted $p$ are in Table~\ref{tab:paired}). A diagnostic on 2,000 SciFact passages points to a compressed embedding space: the effective rank (the exponential of the entropy of the normalised squared singular values of the centred embeddings) falls from 138.5 for the base model to 27.5 for the published 3-fold model (Table~\ref{tab:aniso}). Mean pairwise cosine rises from 0.129 to 0.517, but that number is not comparable across models (BGE-small has 0.628 and retrieves well), so we rely on effective rank. We did not establish the cause; pseudo-labelled in-corpus training is one candidate.

\begin{table}[h]\centering\small
\begin{tabular}{lcc}\toprule
Model & mean pair cosine & effective rank\\\midrule
Base MiniLM & 0.129 & 138.5\\
Published 3-fold & 0.517 & 27.5\\
Ours, averaged & 0.123 & 137.6\\
BGE-small (reference) & 0.628 & 162.5\\\bottomrule\end{tabular}
\caption{Geometry diagnostic on 2,000 random SciFact passages (\texttt{paper/results/anisotropy.json}).}\label{tab:aniso}
\end{table}

\paragraph{Our re-trained models.} They do not show compression, but their nDCG@10 is below the base model's on all four benchmarks: by 0.065 on SciFact, 0.023 on NFCorpus, 0.029 on TREC-COVID and 0.095 on MedQuAD. After Holm correction the differences are clear on SciFact, NFCorpus and MedQuAD; on TREC-COVID (50 queries) the interval for the difference includes zero ($-0.067$ to $+0.007$). Weight averaging over the three seeds changes little: the seed-to-seed range is much smaller than the gap to the base model. Fusing with BM25 raises our model's score on every set, but much of that comes from BM25 itself (BM25 alone scores 0.652 on SciFact against 0.656 for the fusion), and the fusion does not reach the strongest external baselines.

\paragraph{The stronger recipe.} On the test split the averaged R2 model is statistically indistinguishable from the base model on SciFact ($-0.006$), NFCorpus ($-0.005$) and TREC-COVID ($-0.017$), and 0.023 below it on MedQuAD (Holm $p<0.001$; Table~\ref{tab:v2}). It is clearly better than the first recipe (MedQuAD test: $-0.105$ against base) but does not exceed the base model anywhere, and it is far below BGE-small on TREC-COVID (0.455 against 0.756).

\begin{table}[t]\centering\small
\begin{tabular}{lcccc}
\toprule
System & SciFact & NFCorpus & TREC-COVID & MedQuAD (test) \\
\midrule
Base MiniLM & 0.645 & 0.316 & 0.472 & 0.524 \\
First recipe (avg.) & 0.580 & 0.293 & 0.443 & 0.419 \\
Stronger recipe R2 (avg.) & 0.639 & 0.311 & 0.455 & 0.501 \\
BGE-small & 0.713 & 0.343 & 0.756 & 0.587 \\
\bottomrule
\end{tabular}

\caption{Test-split nDCG@10 for the base model, the first recipe, the stronger recipe (both averaged over three seeds) and BGE-small. MedQuAD uses the six test sources (4,599 queries).}\label{tab:v2}
\end{table}

\paragraph{External baselines.} BGE-small and GTE-small beat the base model on all four sets. E5-small beats it on TREC-COVID and MedQuAD but not clearly after correction on SciFact and NFCorpus. MedCPT beats it on SciFact and NFCorpus, is not clearly different on TREC-COVID, and is below it on MedQuAD (0.535 against 0.642). We did not investigate the MedQuAD result; its query encoder was trained for PubMed-style queries, and here the task is consumer questions against long answers truncated to the model's maximum length. BM25 is competitive on SciFact and TREC-COVID and weak on MedQuAD, where questions and answers share little vocabulary.

\begin{table*}[t]\centering\small
\begin{tabular}{llccc}
\toprule
Dataset & System (vs. base) & $\Delta$ nDCG@10 & 95\% CI & Holm $p$ \\
\midrule
SciFact & BGE-small & +0.068 & [+0.041, +0.094] & 0.0046 \\
SciFact & Published HealthMate 3-fold & -0.396 & [-0.442, -0.348] & 0.0046 \\
SciFact & MedCPT & +0.090 & [+0.057, +0.124] & 0.0046 \\
SciFact & E5-small & +0.042 & [+0.014, +0.071] & 0.0343 \\
SciFact & Ours, seed-averaged & -0.065 & [-0.092, -0.036] & 0.0046 \\
SciFact & GTE-small & +0.082 & [+0.054, +0.111] & 0.0046 \\
SciFact & Ours vs. MedCPT & -0.154 & [-0.194, -0.115] & 0.0046 \\
\midrule
NFCorpus & Published HealthMate 3-fold & -0.247 & [-0.278, -0.217] & 0.0046 \\
NFCorpus & GTE-small & +0.033 & [+0.018, +0.048] & 0.0046 \\
NFCorpus & Ours, seed-averaged & -0.023 & [-0.035, -0.011] & 0.0046 \\
NFCorpus & E5-small & +0.009 & [-0.006, +0.026] & 0.9668 \\
NFCorpus & MedCPT & +0.037 & [+0.016, +0.060] & 0.0060 \\
NFCorpus & BGE-small & +0.027 & [+0.013, +0.042] & 0.0077 \\
NFCorpus & Ours vs. MedCPT & -0.060 & [-0.085, -0.037] & 0.0046 \\
\midrule
TREC-COVID & E5-small & +0.271 & [+0.199, +0.345] & 0.0046 \\
TREC-COVID & BGE-small & +0.284 & [+0.212, +0.359] & 0.0046 \\
TREC-COVID & Published HealthMate 3-fold & -0.253 & [-0.334, -0.174] & 0.0046 \\
\midrule
MedQuAD & E5-small & +0.071 & [+0.067, +0.075] & 0.0046 \\
MedQuAD & GTE-small & +0.083 & [+0.079, +0.086] & 0.0046 \\
MedQuAD & MedCPT & -0.108 & [-0.114, -0.101] & 0.0046 \\
MedQuAD & BGE-small & +0.081 & [+0.077, +0.085] & 0.0046 \\
MedQuAD & Ours, seed-averaged & -0.095 & [-0.099, -0.092] & 0.0046 \\
MedQuAD & Published HealthMate 3-fold & -0.463 & [-0.469, -0.456] & 0.0046 \\
MedQuAD & Ours vs. GTE-small & -0.178 & [-0.182, -0.174] & 0.0046 \\
\bottomrule
\end{tabular}

\caption{Paired bootstrap differences in nDCG@10 against the base model (10,000 resamples), with Holm-adjusted $p$ over all comparisons in the analysis. The last row of each block compares our model with the best external baseline (chosen post hoc, so exploratory).}\label{tab:paired}
\end{table*}

\section{Generation: a small local study}\label{sec:rag}
The project also released LoRA adapters for Mistral-7B-Instruct-v0.2. We checked every tensor in the three published adapter files (seeds 42, 123, 999): all 448 tensors in each file contain NaN (\texttt{paper/results/adapter\_nan.json}), and the three files are byte-identical (same SHA-256, \texttt{paper/results/adapter\_hashes.json}), so the ``three seeds'' are one file, so any run through them returns NaN logits and we report no adapter comparison. The project's training logs, with a mean loss near 12, are consistent with a diverged run, but we did not establish the cause.

We evaluate the base generator on PubMedQA \citep{jin2019pubmedqa} questions with a retrieval-augmented setup in which the answer is findable: the corpus is the 1,000 abstracts of the expert-labelled set, so a gold abstract exists for every question. We use the 442 questions (from a seeded sample of 500) labelled yes or no, put the top-1 retrieved abstract in the prompt, and score the first token ``yes'' against ``no'' with no sampling. Table~\ref{tab:rag2} gives the result. Without context the generator scores 0.631; with the gold abstract it scores 0.803, so the test can detect the effect of retrieval. The base, our averaged, R2 and BGE-small retrievers find the gold abstract for 93\% to 99\% of questions and reach 0.792 to 0.803. The published 3-fold encoder finds it for only 60\% and reaches 0.738, still above no context but below the base retriever by 0.061 (95\% CI $-0.090$ to $-0.032$; exploratory and not adjusted for multiple comparisons). Because the corpus is small and contains the gold abstract, retrieval is easy for every sound encoder, and the test is mostly sensitive to bad retrievers. Always answering ``yes'' scores 0.624. We measured neither faithfulness nor hallucination.

\paragraph{An earlier design that we discarded.} Our first design (Amendment 3) retrieved from MedlinePlus chunks that cannot answer PubMedQA questions, and scored yes, no and maybe. The base generator answered ``maybe'' for about 90\% of questions in every condition (accuracy about 0.16), and retrieval changed accuracy by at most 0.006. We report this only because it was run and preregistered; it is superseded because it measured the prompt's hedging, not retrieval.

\begin{table}[t]\centering\small
\begin{tabular}{lccc}
\toprule
Context & Hit@1 & Accuracy [95\% CI] & $\Delta$ vs. none \\
\midrule
None & -- & 0.631 [0.586, 0.676] & -- \\
Gold abstract & -- & 0.803 [0.767, 0.839] & +0.172 [+0.129, +0.215] \\
Base MiniLM & 0.975 & 0.799 [0.760, 0.835] & +0.167 [+0.124, +0.210] \\
First recipe (avg.) & 0.930 & 0.792 [0.753, 0.830] & +0.161 [+0.120, +0.201] \\
Published 3-fold & 0.604 & 0.738 [0.697, 0.778] & +0.106 [+0.068, +0.147] \\
BGE-small & 0.991 & 0.803 [0.765, 0.839] & +0.172 [+0.129, +0.215] \\
Stronger recipe R2 (avg.) & 0.966 & 0.799 [0.760, 0.835] & +0.167 [+0.124, +0.210] \\
\bottomrule
\end{tabular}

\caption{PubMedQA yes/no (442 questions), base Mistral-7B-Instruct-v0.2 with top-1 retrieved abstract. Hit@1 = gold abstract retrieved. Difference vs. no context with paired bootstrap CI.}\label{tab:rag2}
\end{table}

\section{Discussion}
The practical reading is limited. For these small encoders and this much open text, self-supervised in-domain fine-tuning did not improve zero-shot retrieval over the base model. The first recipe hurt, the stronger recipe removed most of the harm, and neither beat general-purpose encoders trained on far more data. The published checkpoints were clearly harmed, and their weaker retrieval went with lower accuracy in the question-answering test. It would be wrong to read this as showing that domain adaptation cannot work; we tried two recipes at one scale on one corpus, and our selection data were themselves a benchmark-style task.

\section*{Limitations}
This paper is mainly a retrieval audit. Its generation study is small and local: 442 PubMedQA yes/no questions, a corpus of 1,000 abstracts that always contains the gold abstract, one base generator, next-token scoring only. The published LoRA adapters could not be evaluated because their weights are NaN, and we did not train replacements. It measures neither faithfulness nor hallucination, and it involved no human raters or clinical outcome. The benchmarks are scientific abstracts and NIH-sourced question--answer text, not real consumer queries, and not the encyclopedia the published models were trained on, so our test is out of domain for them. The original training data and code were unavailable, so we audit behaviour, not the training procedure. Our re-training covers two recipes on one small public-domain corpus; the first used a saturated validation set and a grid whose best point was at its edge, and the second used a new development split but a small grid and one replay corpus whose licence we did not confirm. We cannot say that no recipe could help. BM25 uses a Python implementation whose scores differ somewhat from the usual Anserini baseline. MedQuAD answers were truncated to each model's maximum input length (256 tokens for MiniLM variants, 512 for BGE, E5 and MedCPT); raising the MiniLM-family limit to 512 lowers the base model to 0.591 and our model to 0.529 nDCG@10 (published 3-fold: 0.164), so their order and the clear gap of 0.063 are unchanged. This check was added after the main results. The comparison of our model with the best external baseline per dataset is post hoc.

\section*{Ethics Statement}
No human subjects or personal data were involved. Data are public (BEIR sets, MedQuAD under CC BY 4.0, MedlinePlus text in the public domain); we redistribute none of it. The audited project's training text came from a copyrighted encyclopedia that had been committed to its public repository; the author removed it from the repository's history and we did not use it. The models are research artefacts, not medical advice tools, and must not be used for diagnosis or treatment decisions.

\section*{LLM Use Disclosure}
An AI coding assistant (Claude Code) wrote and ran the evaluation code, searched for and verified references, and drafted this text under the author's direction. The author is responsible for every claim and must confirm each against the repository before submission. Venue-specific disclosure wording must be checked against the final call for papers.

\bibliography{references}
\bibliographystyle{acl_natbib}
\end{document}
